% ===========================================================================
% SECTION 8 — CONCLUSION
% Target: 300-400 words
% ===========================================================================

\section{Conclusion}
\label{sec:conclusion}

This paper presented a systematic comparison of six machine learning
surrogate architectures — LLNFM, LSTM, GP, PINN, TCN, and SW-MLP, in both
a V1 and an improved V2 configuration — integrated as internal state
predictors within a nonlinear MPC framework for rotor speed regulation of
the NREL 5-MW reference wind turbine in above-rated operation. Four
contributions follow from the results reported in
Sections~\ref{sec:ml_surrogates}--\ref{sec:closed_loop_results}:
(1) a quantitative demonstration that one-step prediction accuracy and
closed-loop MPC performance are only loosely coupled, with the most
accurate open-loop surrogates (LLNFM, GP-v2) not translating that
advantage directly into the best closed-loop tracking; (2) a five-stage
diagnostic (Section~\ref{subsec:operational_feasibility}) tracing an
apparent hard operational incompatibility between LSTM, TCN, SW-MLP,
PINN-v2, and GP-v2 on one hand and SQP-based \texttt{nlmpc} on the
other to a specific, resolvable cause — repeated \texttt{predict()}
calls inside the solver's iteration loop, not surrogate architecture,
model size, or MATLAB version — together with a verified manual
forward-pass remedy restoring closed-loop viability for all six
surrogates at speedups of $7.8\times$ to $2844\times$
(Table~\ref{tab:predict_bypass}); (3) a Region~II/III generator torque
switching correction, calibrated at the turbine's actual rated operating
point rather than its theoretical aerodynamic optimum, without which
long-horizon closed-loop simulation enters an irrecoverable trap; and
(4) a Gaussian-Process-based uncertainty-penalized MPC cost, made
statistically meaningful by Matérn~5/2 kernel selection and Bayesian
hyperparameter optimization, which achieves zero pitch activity while
matching the best closed-loop tracking accuracy among the controllers
evaluated in the original (pre-bypass) protocol.

For practitioners selecting an ML surrogate for wind turbine MPC, our
revised results temper the earlier, narrower recommendation to restrict
consideration to instantaneous-state surrogates (LLNFM, GP): once
\texttt{predict()}'s call overhead is bypassed with a verified manual
forward pass, all six architectures tested here — including TCN and
LSTM, previously the most computationally prohibitive — become viable
within the $T_s=\SI{0.1}{s}$ real-time budget. The practical decision
rule is accordingly not architectural but methodological: before
concluding that a given surrogate is operationally infeasible inside an
SQP-based \texttt{nlmpc} solve, practitioners should profile whether
the surrogate's own inference call, or MATLAB's \texttt{predict()}
dispatch overhead specifically, is the dominant cost, since only the
former is an intrinsic property of the architecture. Among the six
surrogates evaluated with the manual bypass
(Section~\ref{subsec:bypass_closed_loop}), LSTM achieved the lowest
mean tracking RMSE at two of three tested wind speeds, at the highest
--- though still real-time-compliant --- per-step cost, reflecting its
irreducible sequential recurrence; the single-pass architectures
(LLNFM, GP, the residual MLP, SW-MLP, PINN-v2) remain preferable where
computational headroom is scarcer.

The main limitations of this study — a two-state plant model, uniform
wind input, an unresolved seed-to-seed instability in LSTM's tracking
performance near the Region~II/III boundary, and the per-architecture
verification burden inherent to any manual-bypass approach — are
addressed in Section~\ref{subsec:limitations}, and motivate the
extensions outlined in Section~\ref{subsec:perspectives_v3}: targeted
investigation of the wind-speed-dependent crossover between surrogate
families and of LSTM's low-wind-speed instability, a four-state plant
model including drivetrain torsion and tower fore-aft dynamics, and an
explicit wind-speed forecasting or estimation layer to complement the
uncertainty-aware MPC cost introduced in this work.
